
%%%%%%%%%%%%%%%%%%%%%%%%%%%%%%%%%%%%%%
%%            Introdução            %%
%%%%%%%%%%%%%%%%%%%%%%%%%%%%%%%%%%%%%%

\chapter{Introdução}
  \label{cha:intro}

  Para \cite{tarapanoff1995tecnicas}, a tomada de decisões é um processo que envolve a escolha de uma ou mais alternativas, a partir de um conjunto de opções, para atingir um objetivo. Este processo é realizado por indivíduos, empresas e governos, e é fundamental para o desenvolvimento de qualquer atividade.

  Para que uma decisão seja tomada, é necessário que o indivíduo ou grupo de indivíduos tenha conhecimento sobre o problema a ser resolvido, bem como sobre as alternativas que podem ser escolhidas para resolvê-lo. O conhecimento sobre o problema e as alternativas é obtido por informações, que podem ser alcançadas atráves do processamento de diversas fontes de dados, como sensores, gravações, registros, entre outros.

  A análise dos dados é realizada através de técnicas e ferramentas, que são utilizadas para extrair conhecimento a partir de dados \cite{guimaraes2004sistema}. A partir do conhecimento obtido, o indivíduo ou grupo de indivíduos deve analisar as informações e selecionar as que são mais relevantes para a tomada de decisão.

  % Contextualização e problema
  \section{Contextualização e Problema}
  \label{subsec:contextualizacao-problema}

  A análise de dados é um dos principais fatores para o sucesso de uma empresa, pois é através dela que é possível obter informações que se tornam valiosas para a tomada de decisões. Logo, dado o aumento da quantidade de dados gerados e a necessidade de se obter informações a partir deles, gerou-se o interesse por ferramentas que auxiliem na análise e visualização de dados.

  Um grande conjunto de ferramentas já está disponível online, porém, a maioria delas não é gratuita, performática e nem de código aberto. A ausência de ferramentas deste tipo é um fator notável, visto que nem sempre as empresas possuem recursos ou pessoal para desenvolver ferramentas deste tipo.

  Nesta perspectiva, surge a oportunidade de se desenvolver uma ferramenta que atenda aos critérios citados anteriormente, e que seja facilmente extensível para novas funcionalidades e diferentes tipos e fontes de dados.




  % Objetivos (geral e específicos)
  \section{Objetivos}
\label{sec:objetivos}

\subsection{Objetivo Geral}
\label{subsec:objetivo-geral}

O objetivo deste trabalho é o estudo e desenvolvimento de um projeto de software com a finalidade de análise e visualização de dados, que seja modular e facilmente extensível, para atender a vários propósitos e fontes de dados.

\subsection{Objetivos Específicos}
\label{subsec:objetivos-especificos}

\begin{itemize}
    \item Identificar diferentes problemas de análise de dados, tipos e fontes de dados.
    \item Modelar um projeto de software que atenda as necessidades identificadas, utilizando os conceitos de \ac{es}.
    \item Implementar o software modelado.
    \item Testar e validar o software para com aplicação em laboratório.
    \item Documentar o software para futuras modificações.
\end{itemize}

  % Delimitação do estudo
  
% Nesta seção deve ser apresentado o problema abordado no estudo de forma a especificá-lo dentro do contexto em que está inserido.

\section{Delimitação do Estudo}
    \label{sec:delimitacao-estudo}

    As ferramentas de visualização e análise de dados possuem um papel fundamental na tomada de decisão, sendo assim sua utilização é cada vez mais comum em diversas áreas. No entanto, a maioria das ferramentas disponíveis no mercado são proprietárias, de código fechado ou foram criadas em linguagens pouco utilizadas, o que diminui sua acessibilidade e dificulta o uso. Além disso, a maioria das ferramentas disponíveis são específicas para um tipo de problema, o que dificulta a sua utilização em diversos contextos.

    Diante disso, o presente trabalho tem como objetivo a extensão de uma ferramente de visualização e análise de dados desenvolvida durante o projeto de pesquisa SmartEnergy. O projeto SmartEnergy tinha como objetivo a análise de dados de consumo energético em séries temporais. Logo, esse trabalho busca realizar o estudo da capacidade de extensão da ferramenta para outros contextos e evolução dela para uma arquitetura modular e de uso mais amplo.

    A ferramenta foi desenvolvida em Python, utilizando a biblioteca Matplotlib para a visualização dos dados. A biblioteca Matplotlib é uma biblioteca de visualização de dados 2D em Python, que produz gráficos de qualidade de publicação em uma variedade de formatos impressos e ambientes interativos em plataformas diversas. A interface gráfica do usuário \ac{gui} foi desenvolvida utilizando a biblioteca PyQt que fornece um conjunto de interfaces para a biblioteca Qt. A biblioteca Qt é uma biblioteca de código aberto para desenvolvimento de aplicações multiplataforma.

  % Justificativa
  
% Esta seção trata os motivos pelos quais seu trabalho é relevante

\section{Justificativa}
    \label{sec:justificativa}

    A tarefa de análise de dados é complexa e exige o conhecimento e emprego de diversas técnicas de processamento \cite{gomes2005organizaccao}. Cada problema requer uma abordagem diferente, e a escolha da técnica mais adequada é um desafio, necessitando de uma vasta gama de ferramentas e conhecimentos.

    O processo exploratório da análise de dados requer o emprego de técnicas de visualização de dados, que permitem a compreensão dos dados e a identificação de padrões e tendências. Logo após a visualização dos dados é identificado o problema a ser resolvido, e então a escolha das técnicas que serão utilizadas para a resolução do problema. Frequentemente, a escolha das técnicas é baseada em experiência e intuição, necessitando de uma série de testes e ajustes dos parâmetros para a obtenção de resultados satisfatórios.

    Além disso, a aplicação de determinadas técnica são realizadas em conjunto, e a escolha da ordem de aplicação das técnicas é um desafio, pois a ordem de aplicação das técnicas pode influenciar no resultado \cite{gomes2005organizaccao}. Por exemplo, a aplicação de técnicas de redução de dimensionalidade antes da aplicação de técnicas de agrupamento pode resultar em um agrupamento mais eficiente. Logo, existe um número grande de possibilidades de combinação de técnicas, e a escolha da melhor combinação é um desafio.

    Desta forma, a proposta deste trabalho é a criação de uma plataforma para análise e visualização de dados, que permita a escolha das técnicas, a ordem de aplicação e a visualização dos resultados. A plataforma será desenvolvida através do levantamento de requisitos necessários e o desenvolvimento de um protótipo que atenda aos mesmos.

  % Apresentação do trabalho
  
% Esta seção trata a construção do documento, especificamente a disposição dos capítulos e uma introdução do que está sendo abordado nestes.

\section{Apresentação do Trabalho}
    \label{sec:apresentacao-trabalho}

    No capítulo seguinte, fundamentação teórica, são discutidos assuntos relacionados ao trabalho, como Análise e visualização de dados, Aprendizado de Máquina, Engenharia de Software e as linguagens e bibliotecas utilizadas no desenvolvimento do sistema proposto.

    No terceiro capítulo são apresentados tópicos relacionados a modelagem do sistema, contendo os requisitos, funcionais e não funcionais, diagramas, modelo de dados e a arquitetura modular do sistema.

    O quarto capítulo é dedicado a implementação do sistema, contendo a descrição do processo de desenvolvimento, a arquitetura do sistema, a descrição das tecnologias utilizadas e a descrição do processo de desenvolvimento.

    Adiante, no quinto capítulo são apresentados os resultados obtidos com o sistema, contendo a descrição da interface, exemplos de uso e a avaliação do sistema.

    No sexto capítulo, são apresentadas a conclusão obtida com o presente trabalho, bem como sugestões de trabalhos futuros. Por fim, são listadas as referências bibliográficas utilizadas durante este trabalho, seguidas por um apêndice com matérias relacionados ao trabalho e um anexo com o código-fonte do sistema.


